\documentclass{beamer}
\usetheme{Madrid}
\usecolortheme{default}
\usepackage[utf8]{inputenc}
\usepackage{amsmath}
\usepackage{amsfonts}
\usepackage{amssymb}
\usepackage{tikz}

\title{CRISC Domain D - 100 More Impossible Hard Questions}
\subtitle{Information Technology and Security - Expert Level Part 4}
\author{CRISC Mastery Series}
\date{\today}

\begin{document}
	
	\begin{frame}
		\titlepage
	\end{frame}
	
	\begin{frame}{Warning}
		\begin{alertblock}{Extreme Difficulty}
			Domain D is technical but the exam tests judgment, not trivia.\\
			Every option will look plausible.\\
			Think like a risk practitioner advising on technology, not a technician.
		\end{alertblock}
	\end{frame}
	
	\begin{frame}{Outline}
		\tableofcontents
	\end{frame}
	
	% ============================================================
	\section{1. Enterprise Architecture}
	% ============================================================
	
	\begin{frame}{Q301: Architecture Principle Enforcement}
		\begin{block}{Question}
			A project manager claims an architecture principle does not apply to their project. What is the BEST governance response?
		\end{block}
		\begin{enumerate}[(A)]
			\item Ignore the claim
			\item Require documented analysis and approval through the architecture review or exception process
			\item Proceed without review
			\item Cancel the project
		\end{enumerate}
	\end{frame}
	
	\begin{frame}{Q301: Explanation}
		\begin{block}{Correct Answer: B}
			\textbf{Claims need governance review:}
			\begin{itemize}
				\item Principles apply by default
				\item Exemptions require review
				\item Documentation protects decisions
			\end{itemize}
		\end{block}
		\begin{alertblock}{Critical Trap}
			Principle exemptions need governance, not assertion.
		\end{alertblock}
	\end{frame}
	
	\begin{frame}{Q302: Architecture Roadmap Risk}
		\begin{block}{Question}
			What is the MAIN risk of an outdated EA roadmap?
		\end{block}
		\begin{enumerate}[(A)]
			\item It looks old
			\item It drives investment toward obsolete or misaligned states
			\item It costs more
			\item It is too detailed
		\end{enumerate}
	\end{frame}
	
	\begin{frame}{Q302: Explanation}
		\begin{block}{Correct Answer: B}
			\textbf{Stale roadmaps misalign investment:}
			\begin{itemize}
				\item Guides investment
				\item Outdated guides wrongly
				\item Review regularly
			\end{itemize}
		\end{block}
		\begin{alertblock}{Critical Trap}
			Stale roadmaps cause misallocation.
		\end{alertblock}
	\end{frame}
	
	\begin{frame}{Q303: Business Capability Mapping}
		\begin{block}{Question}
			What is the purpose of business capability mapping?
		\end{block}
		\begin{enumerate}[(A)]
			\item Marketing
			\item Describe what the business can do independent of organizational structure
			\item Reduce staff
			\item Buy tools
		\end{enumerate}
	\end{frame}
	
	\begin{frame}{Q303: Explanation}
		\begin{block}{Correct Answer: B}
			\textbf{Capabilities describe what business can do:}
			\begin{itemize}
				\item Business outcomes
				\item Stable view
				\item Independent of structure
			\end{itemize}
		\end{block}
		\begin{alertblock}{Critical Trap}
			Capabilities describe business outcomes, not org charts.
		\end{alertblock}
	\end{frame}
	
	\begin{frame}{Q304: Architecture Compliance Review}
		\begin{block}{Question}
			What should an architecture compliance review verify?
		\end{block}
		\begin{enumerate}[(A)]
			\item Only cost
			\item Alignment with principles, standards, and target architecture
			\item Only vendor logos
			\item Only timelines
		\end{enumerate}
	\end{frame}
	
	\begin{frame}{Q304: Explanation}
		\begin{block}{Correct Answer: B}
			\textbf{Compliance means alignment:}
			\begin{itemize}
				\item Principles
				\item Standards
				\item Target architecture
			\end{itemize}
		\end{block}
		\begin{alertblock}{Critical Trap}
			Compliance review is alignment, not cost check.
		\end{alertblock}
	\end{frame}
	
	\begin{frame}{Q305: Application Rationalization}
		\begin{block}{Question}
			What is the main benefit of application rationalization?
		\end{block}
		\begin{enumerate}[(A)]
			\item More apps
			\item Reduced duplication, cost, and complexity
			\item Higher licensing
			\item Better visuals
		\end{enumerate}
	\end{frame}
	
	\begin{frame}{Q305: Explanation}
		\begin{block}{Correct Answer: B}
			\textbf{Rationalization reduces duplication:}
			\begin{itemize}
				\item Consolidate overlapping apps
				\item Retire low-value systems
				\item Lower cost and complexity
			\end{itemize}
		\end{block}
		\begin{alertblock}{Critical Trap}
			Rationalization reduces cost and complexity.
		\end{alertblock}
	\end{frame}
	
	\begin{frame}{Q306: Architecture Principle Characteristics}
		\begin{block}{Question}
			Which is a characteristic of a good architecture principle?
		\end{block}
		\begin{enumerate}[(A)]
			\item Vague
			\item Understandable, relevant, and actionable
			\item Extreme
			\item Optional
		\end{enumerate}
	\end{frame}
	
	\begin{frame}{Q306: Explanation}
		\begin{block}{Correct Answer: B}
			\textbf{Good principles are clear and actionable:}
			\begin{itemize}
				\item Understandable
				\item Relevant
				\item Actionable
			\end{itemize}
		\end{block}
		\begin{alertblock}{Critical Trap}
			Vague principles are unusable.
		\end{alertblock}
	\end{frame}
	
	\begin{frame}{Q307: Architecture Debt Visibility}
		\begin{block}{Question}
			Why should architecture debt be visible to leadership?
		\end{block}
		\begin{enumerate}[(A)]
			\item It looks bad
			\item It affects cost, risk, and speed of change
			\item It is entertainment
			\item It is optional
		\end{enumerate}
	\end{frame}
	
	\begin{frame}{Q307: Explanation}
		\begin{block}{Correct Answer: B}
			\textbf{Debt affects business outcomes:}
			\begin{itemize}
				\item Costs
				\item Risk
				\item Change speed
			\end{itemize}
		\end{block}
		\begin{alertblock}{Critical Trap}
			Hidden debt becomes a business surprise.
		\end{alertblock}
	\end{frame}
	
	\begin{frame}{Q308: EA and Merger Integration}
		\begin{block}{Question}
			What is the role of EA in a merger integration?
		\end{block}
		\begin{enumerate}[(A)]
			\item Ignore
			\item Guide consolidated target architecture and integration priorities
			\item Marketing
			\item Payroll
		\end{enumerate}
	\end{frame}
	
	\begin{frame}{Q308: Explanation}
		\begin{block}{Correct Answer: B}
			\textbf{EA guides integration:}
			\begin{itemize}
				\item Consolidate architectures
				\item Sequence priorities
				\item Reduce duplication
			\end{itemize}
		\end{block}
		\begin{alertblock}{Critical Trap}
			M\&A integration needs EA direction.
		\end{alertblock}
	\end{frame}
	
	\begin{frame}{Q309: Architecture Governance Escalation}
		\begin{block}{Question}
			If an architecture decision cannot be resolved at the working level, what should happen?
		\end{block}
		\begin{enumerate}[(A)]
			\item Ignore it
			\item Escalate to appropriate governance authority
			\item Pick randomly
			\item Cancel
		\end{enumerate}
	\end{frame}
	
	\begin{frame}{Q309: Explanation}
		\begin{block}{Correct Answer: B}
			\textbf{Escalation is a control:}
			\begin{itemize}
				\item Defined escalation path
				\item Authority resolves
				\item Documented outcome
			\end{itemize}
		\end{block}
		\begin{alertblock}{Critical Trap}
			Unresolved decisions stall projects.
		\end{alertblock}
	\end{frame}
	
	\begin{frame}{Q310: EA Value Communication}
		\begin{block}{Question}
			How should EA value be communicated to business leaders?
		\end{block}
		\begin{enumerate}[(A)]
			\item Technical diagrams only
			\item Business outcomes enabled, cost savings, and reduced risk
			\item Vendor names
			\item Only budgets
		\end{enumerate}
	\end{frame}
	
	\begin{frame}{Q310: Explanation}
		\begin{block}{Correct Answer: B}
			\textbf{Speak business outcomes:}
			\begin{itemize}
				\item Outcomes
				\item Savings
				\item Risk reduction
			\end{itemize}
		\end{block}
		\begin{alertblock}{Critical Trap}
			Business leaders want outcomes, not diagrams.
		\end{alertblock}
	\end{frame}
	
	% ============================================================
	\section{2. IT Operations and Change Management}
	% ============================================================
	
	\begin{frame}{Q311: Change Risk Assessment}
		\begin{block}{Question}
			What should a change risk assessment evaluate?
		\end{block}
		\begin{enumerate}[(A)]
			\item Only cost
			\item Business impact, likelihood of failure, and dependencies
			\item Only vendor names
			\item Only deadlines
		\end{enumerate}
	\end{frame}
	
	\begin{frame}{Q311: Explanation}
		\begin{block}{Correct Answer: B}
			\textbf{Assess impact and likelihood:}
			\begin{itemize}
				\item Business impact
				\item Failure likelihood
				\item Dependencies
			\end{itemize}
		\end{block}
		\begin{alertblock}{Critical Trap}
			Change risk = impact times likelihood.
		\end{alertblock}
	\end{frame}
	
	\begin{frame}{Q312: Standard Change Model}
		\begin{block}{Question}
			What is a standard change?
		\end{block}
		\begin{enumerate}[(A)]
			\item A new emergency
			\item A pre-approved, low-risk, repeatable change
			\item A vendor change
			\item A cancelled change
		\end{enumerate}
	\end{frame}
	
	\begin{frame}{Q312: Explanation}
		\begin{block}{Correct Answer: B}
			\textbf{Standard changes are pre-approved:}
			\begin{itemize}
				\item Pre-approved
				\item Low risk
				\item Repeatable
			\end{itemize}
		\end{block}
		\begin{alertblock}{Critical Trap}
			Standard changes skip per-change approval.
		\end{alertblock}
	\end{frame}
	
	\begin{frame}{Q313: CAB Emergency Path}
		\begin{block}{Question}
			What CAB mechanism supports urgent changes?
		\end{block}
		\begin{enumerate}[(A)]
			\item Ignore
			\item Emergency change process with retrospective review
			\item Remove CAB
			\item Proceed without records
		\end{enumerate}
	\end{frame}
	
	\begin{frame}{Q313: Explanation}
		\begin{block}{Correct Answer: B}
			\textbf{Emergency with retrospective:}
			\begin{itemize}
				\item Emergency path
				\item Retrospective review
				\item Documented
			\end{itemize}
		\end{block}
		\begin{alertblock}{Critical Trap}
			Emergency still requires governance.
		\end{alertblock}
	\end{frame}
	
	\begin{frame}{Q314: Release Risk}
		\begin{block}{Question}
			Which factor MOST affects release risk?
		\end{block}
		\begin{enumerate}[(A)]
			\item Release time
			\item Number and complexity of changes in the release
			\item Vendor logo
			\item Marketing
		\end{enumerate}
	\end{frame}
	
	\begin{frame}{Q314: Explanation}
		\begin{block}{Correct Answer: B}
			\textbf{Release risk from volume and complexity:}
			\begin{itemize}
				\item Volume of changes
				\item Complexity
				\item Dependencies
			\end{itemize}
		\end{block}
		\begin{alertblock}{Critical Trap}
			Release risk rises with change volume and complexity.
		\end{alertblock}
	\end{frame}
	
	\begin{frame}{Q315: Incident Priority}
		\begin{block}{Question}
			Which factor drives incident priority in ITIL?
		\end{block}
		\begin{enumerate}[(A)]
			\item Marketing
			\item Impact and urgency combined
			\item Number of meetings
			\item Vendor discounts
		\end{enumerate}
	\end{frame}
	
	\begin{frame}{Q315: Explanation}
		\begin{block}{Correct Answer: B}
			\textbf{Priority from impact and urgency:}
			\begin{itemize}
				\item Impact
				\item Urgency
				\item Combined priority
			\end{itemize}
		\end{block}
		\begin{alertblock}{Critical Trap}
			Priority = impact plus urgency.
		\end{alertblock}
	\end{frame}
	
	\begin{frame}{Q316: Problem vs Incident}
		\begin{block}{Question}
			What is the main output of problem management?
		\end{block}
		\begin{enumerate}[(A)]
			\item Faster restore only
			\item Known errors and permanent fixes to remove root causes
			\item New tickets
			\item More vendors
		\end{enumerate}
	\end{frame}
	
	\begin{frame}{Q316: Explanation}
		\begin{block}{Correct Answer: B}
			\textbf{Problem management reduces recurrence:}
			\begin{itemize}
				\item Root causes
				\item Known errors
				\item Permanent fixes
			\end{itemize}
		\end{block}
		\begin{alertblock}{Critical Trap}
			Problem management eliminates recurrence.
		\end{alertblock}
	\end{frame}
	
	\begin{frame}{Q317: Knowledge Management Value}
		\begin{block}{Question}
			Why is knowledge management valuable in IT operations?
		\end{block}
		\begin{enumerate}[(A)]
			\item Marketing
			\item Reduces resolution time and improves consistency by documenting solutions
			\item Reduces cost only
			\item Removes staff
		\end{enumerate}
	\end{frame}
	
	\begin{frame}{Q317: Explanation}
		\begin{block}{Correct Answer: B}
			\textbf{Knowledge management improves resolution:}
			\begin{itemize}
				\item Documented solutions
				\item Faster resolution
				\item Consistent handling
			\end{itemize}
		\end{block}
		\begin{alertblock}{Critical Trap}
			Knowledge management scales operations.
		\end{alertblock}
	\end{frame}
	
	\begin{frame}{Q318: Service Catalog Purpose}
		\begin{block}{Question}
			What is the purpose of a service catalog?
		\end{block}
		\begin{enumerate}[(A)]
			\item Marketing
			\item A structured list of services offered with descriptions and SLAs
			\item Employee records
			\item Budgets
		\end{enumerate}
	\end{frame}
	
	\begin{frame}{Q318: Explanation}
		\begin{block}{Correct Answer: B}
			\textbf{Service catalog lists services:}
			\begin{itemize}
				\item Structured list
				\item Descriptions
				\item SLAs
			\end{itemize}
		\end{block}
		\begin{alertblock}{Critical Trap}
			Service catalog is the front door to IT.
		\end{alertblock}
	\end{frame}
	
	\begin{frame}{Q319: Capacity Management}
		\begin{block}{Question}
			What is the purpose of capacity management?
		\end{block}
		\begin{enumerate}[(A)]
			\item Marketing
			\item Ensure IT resources meet current and future demand at acceptable cost
			\item Reduce headcount
			\item Increase cost
		\end{enumerate}
	\end{frame}
	
	\begin{frame}{Q319: Explanation}
		\begin{block}{Correct Answer: B}
			\textbf{Capacity meets demand:}
			\begin{itemize}
				\item Current demand
				\item Future demand
				\item Acceptable cost
			\end{itemize}
		\end{block}
		\begin{alertblock}{Critical Trap}
			Capacity management prevents surprises.
		\end{alertblock}
	\end{frame}
	
	\begin{frame}{Q320: Availability Management}
		\begin{block}{Question}
			What does availability management aim to deliver?
		\end{block}
		\begin{enumerate}[(A)]
			\item Marketing
			\item Agreed availability targets aligned to business needs
			\item Higher cost
			\item Fewer tickets
		\end{enumerate}
	\end{frame}
	
	\begin{frame}{Q320: Explanation}
		\begin{block}{Correct Answer: B}
			\textbf{Availability matches business needs:}
			\begin{itemize}
				\item Availability targets
				\item Business alignment
				\item Measured
			\end{itemize}
		\end{block}
		\begin{alertblock}{Critical Trap}
			Availability is about meeting business expectations.
		\end{alertblock}
	\end{frame}
	
	% ============================================================
	\section{3. Enterprise Resiliency}
	% ============================================================
	
	\begin{frame}{Q321: DR Plan Scope}
		\begin{block}{Question}
			What should a disaster recovery plan define?
		\end{block}
		\begin{enumerate}[(A)]
			\item Only backups
			\item Roles, procedures, recovery steps, and communication for IT recovery
			\item Only vendor names
			\item Only budgets
		\end{enumerate}
	\end{frame}
	
	\begin{frame}{Q321: Explanation}
		\begin{block}{Correct Answer: B}
			\textbf{DR plan defines recovery:}
			\begin{itemize}
				\item Roles
				\item Procedures
				\item Steps
				\item Communication
			\end{itemize}
		\end{block}
		\begin{alertblock}{Critical Trap}
			DR plan focuses on IT recovery, not just backups.
		\end{alertblock}
	\end{frame}
	
	\begin{frame}{Q322: BCP vs DR Scope}
		\begin{block}{Question}
			Which statement is correct?
		\end{block}
		\begin{enumerate}[(A)]
			\item BCP and DR are the same
			\item BCP covers the business broadly; DR covers IT recovery
			\item DR covers the business broadly
			\item BCP only covers servers
		\end{enumerate}
	\end{frame}
	
	\begin{frame}{Q322: Explanation}
		\begin{block}{Correct Answer: B}
			\textbf{BCP broader than DR:}
			\begin{itemize}
				\item BCP is business-wide
				\item DR is IT-focused
				\item DR supports BCP
			\end{itemize}
		\end{block}
		\begin{alertblock}{Critical Trap}
			BCP is business-wide; DR is one component.
		\end{alertblock}
	\end{frame}
	
	\begin{frame}{Q323: Crisis Communication Plan}
		\begin{block}{Question}
			Why is a crisis communication plan essential?
		\end{block}
		\begin{enumerate}[(A)]
			\item Marketing
			\item Provides coordinated and timely communication to stakeholders during disruption
			\item Reduces cost
			\item Avoids audits
		\end{enumerate}
	\end{frame}
	
	\begin{frame}{Q323: Explanation}
		\begin{block}{Correct Answer: B}
			\textbf{Communication under crisis:}
			\begin{itemize}
				\item Coordinated
				\item Timely
				\item Stakeholders informed
			\end{itemize}
		\end{block}
		\begin{alertblock}{Critical Trap}
			Crisis communication protects trust.
		\end{alertblock}
	\end{frame}
	
	\begin{frame}{Q324: Third-Party BCP Dependency}
		\begin{block}{Question}
			Why should third-party BCP dependencies be assessed?
		\end{block}
		\begin{enumerate}[(A)]
			\item Trust only
			\item Verify providers can meet recovery needs aligned to your RTO and RPO
			\item Cheaper
			\item Faster
		\end{enumerate}
	\end{frame}
	
	\begin{frame}{Q324: Explanation}
		\begin{block}{Correct Answer: B}
			\textbf{Verify alignment:}
			\begin{itemize}
				\item RTO alignment
				\item RPO alignment
				\item Contracts plus verification
			\end{itemize}
		\end{block}
		\begin{alertblock}{Critical Trap}
			Contracts do not guarantee recovery.
		\end{alertblock}
	\end{frame}
	
	\begin{frame}{Q325: DR Exercise Types}
		\begin{block}{Question}
			Which exercise type provides the HIGHEST assurance?
		\end{block}
		\begin{enumerate}[(A)]
			\item Tabletop
			\item Full failover with real workloads
			\item Document review
			\item Walkthrough only
		\end{enumerate}
	\end{frame}
	
	\begin{frame}{Q325: Explanation}
		\begin{block}{Correct Answer: B}
			\textbf{Full failover = highest assurance:}
			\begin{itemize}
				\item Real failover
				\item Real workloads
				\item Higher cost
			\end{itemize}
		\end{block}
		\begin{alertblock}{Critical Trap}
			Full failover tests real recovery.
		\end{alertblock}
	\end{frame}
	
	\begin{frame}{Q326: Backup Testing Frequency}
		\begin{block}{Question}
			Why should backup restores be tested regularly?
		\end{block}
		\begin{enumerate}[(A)]
			\item Cost only
			\item Confirm recovery works and backups are not corrupted
			\item To satisfy marketing
			\item Ignore
		\end{enumerate}
	\end{frame}
	
	\begin{frame}{Q326: Explanation}
		\begin{block}{Correct Answer: B}
			\textbf{Restore tests confirm capability:}
			\begin{itemize}
				\item Recovery works
				\item Data is intact
				\item Cadence agreed
			\end{itemize}
		\end{block}
		\begin{alertblock}{Critical Trap}
			Backups without restore testing are unreliable.
		\end{alertblock}
	\end{frame}
	
	\begin{frame}{Q327: Cloud DR Trade-offs}
		\begin{block}{Question}
			What is the main trade-off in cloud DR?
		\end{block}
		\begin{enumerate}[(A)]
			\item Free vs paid
			\item Cost vs recovery speed and capacity
			\item Only latency
			\item Only firewalls
		\end{enumerate}
	\end{frame}
	
	\begin{frame}{Q327: Explanation}
		\begin{block}{Correct Answer: B}
			\textbf{Cost vs speed:}
			\begin{itemize}
				\item Lower RTO means higher cost
				\item Capacity trade-offs
				\item Match to business
			\end{itemize}
		\end{block}
		\begin{alertblock}{Critical Trap}
			Faster recovery costs more.
		\end{alertblock}
	\end{frame}
	
	\begin{frame}{Q328: BC Plan Ownership}
		\begin{block}{Question}
			Who should own the business continuity plan?
		\end{block}
		\begin{enumerate}[(A)]
			\item IT only
			\item Business leaders with IT support
			\item Vendor
			\item Auditor
		\end{enumerate}
	\end{frame}
	
	\begin{frame}{Q328: Explanation}
		\begin{block}{Correct Answer: B}
			\textbf{Business owns BCP:}
			\begin{itemize}
				\item Business owners
				\item IT supports
				\item Accountability aligned
			\end{itemize}
		\end{block}
		\begin{alertblock}{Critical Trap}
			BCP is a business plan, not an IT plan.
		\end{alertblock}
	\end{frame}
	
	\begin{frame}{Q329: Crisis Decision Rights}
		\begin{block}{Question}
			What is the best practice for crisis decision rights?
		\end{block}
		\begin{enumerate}[(A)]
			\item No rights
			\item Predefined crisis roles, thresholds, and escalation paths
			\item Everything to the board
			\item Random
		\end{enumerate}
	\end{frame}
	
	\begin{frame}{Q329: Explanation}
		\begin{block}{Correct Answer: B}
			\textbf{Predefine crisis rights:}
			\begin{itemize}
				\item Predefined roles
				\item Thresholds
				\item Escalation paths
			\end{itemize}
		\end{block}
		\begin{alertblock}{Critical Trap}
			Clarity enables timely crisis action.
		\end{alertblock}
	\end{frame}
	
	\begin{frame}{Q330: Recovery Objective Alignment}
		\begin{block}{Question}
			Why must RTO and RPO be agreed with business owners?
		\end{block}
		\begin{enumerate}[(A)]
			\item Marketing
			\item Owners understand business impact and acceptable loss
			\item Cheaper
			\item Faster
		\end{enumerate}
	\end{frame}
	
	\begin{frame}{Q330: Explanation}
		\begin{block}{Correct Answer: B}
			\textbf{Owners define impact:}
			\begin{itemize}
				\item Business impact
				\item Acceptable loss
				\item Owner-defined objectives
			\end{itemize}
		\end{block}
		\begin{alertblock}{Critical Trap}
			RTO and RPO are business decisions.
		\end{alertblock}
	\end{frame}
	
	% ============================================================
	\section{4. Data Lifecycle Management}
	% ============================================================
	
	\begin{frame}{Q331: Data Classification Owner}
		\begin{block}{Question}
			Who classifies data?
		\end{block}
		\begin{enumerate}[(A)]
			\item Auditor
			\item Data owner or steward
			\item Vendor
			\item User
		\end{enumerate}
	\end{frame}
	
	\begin{frame}{Q331: Explanation}
		\begin{block}{Correct Answer: B}
			\textbf{Owners classify:}
			\begin{itemize}
				\item Owners classify
				\item Custodians enforce
				\item Auditors verify
			\end{itemize}
		\end{block}
		\begin{alertblock}{Critical Trap}
			Owners classify; custodians implement.
		\end{alertblock}
	\end{frame}
	
	\begin{frame}{Q332: Data Labeling Purpose}
		\begin{block}{Question}
			Why is labeling data important?
		\end{block}
		\begin{enumerate}[(A)]
			\item Marketing
			\item Enables correct handling and controls based on sensitivity
			\item Slower
			\item Bigger files
		\end{enumerate}
	\end{frame}
	
	\begin{frame}{Q332: Explanation}
		\begin{block}{Correct Answer: B}
			\textbf{Labels enable controls:}
			\begin{itemize}
				\item Makes classification actionable
				\item Supports DLP
				\item Guides handling
			\end{itemize}
		\end{block}
		\begin{alertblock}{Critical Trap}
			Labels make classification operational.
		\end{alertblock}
	\end{frame}
	
	\begin{frame}{Q333: Retention Legal Conflict}
		\begin{block}{Question}
			What should happen if a retention schedule conflicts with a regulator's request?
		\end{block}
		\begin{enumerate}[(A)]
			\item Destroy data
			\item Follow legal guidance and preserve as required
			\item Ignore
			\item Archive indefinitely
		\end{enumerate}
	\end{frame}
	
	\begin{frame}{Q333: Explanation}
		\begin{block}{Correct Answer: B}
			\textbf{Legal obligations over schedule:}
			\begin{itemize}
				\item Legal holds or requests
				\item Preserve
				\item Document rationale
			\end{itemize}
		\end{block}
		\begin{alertblock}{Critical Trap}
			Regulatory requests override retention schedules.
		\end{alertblock}
	\end{frame}
	
	\begin{frame}{Q334: Data Destruction Standards}
		\begin{block}{Question}
			Which standards best support secure destruction?
		\end{block}
		\begin{enumerate}[(A)]
			\item Marketing
			\item NIST 800-88, ISO 21964, and industry standards
			\item Only deleting files
			\item Only formatting
		\end{enumerate}
	\end{frame}
	
	\begin{frame}{Q334: Explanation}
		\begin{block}{Correct Answer: B}
			\textbf{Use standards:}
			\begin{itemize}
				\item NIST 800-88
				\item ISO 21964
				\item Industry methods
			\end{itemize}
		\end{block}
		\begin{alertblock}{Critical Trap}
			Deleting or formatting is not destruction.
		\end{alertblock}
	\end{frame}
	
	\begin{frame}{Q335: Data Minimization Benefit}
		\begin{block}{Question}
			What is the main benefit of data minimization?
		\end{block}
		\begin{enumerate}[(A)]
			\item Marketing
			\item Reduces exposure, cost, and compliance burden
			\item Faster
			\item Slower
		\end{enumerate}
	\end{frame}
	
	\begin{frame}{Q335: Explanation}
		\begin{block}{Correct Answer: B}
			\textbf{Less data, less risk:}
			\begin{itemize}
				\item Reduced exposure
				\item Lower cost
				\item Simpler compliance
			\end{itemize}
		\end{block}
		\begin{alertblock}{Critical Trap}
			Data minimization reduces risk and cost.
		\end{alertblock}
	\end{frame}
	
	\begin{frame}{Q336: DSR Response Accuracy}
		\begin{block}{Question}
			What should accompany a response to a data subject access request?
		\end{block}
		\begin{enumerate}[(A)]
			\item Only data
			\item Clear information about the processing and rights
			\item Marketing
			\item Nothing
		\end{enumerate}
	\end{frame}
	
	\begin{frame}{Q336: Explanation}
		\begin{block}{Correct Answer: B}
			\textbf{Include context:}
			\begin{itemize}
				\item What was processed
				\item Purpose
				\item Rights
			\end{itemize}
		\end{block}
		\begin{alertblock}{Critical Trap}
			DSAR needs transparent explanation.
		\end{alertblock}
	\end{frame}
	
	\begin{frame}{Q337: Data Residency in Backup}
		\begin{block}{Question}
			Why do backups need data residency alignment?
		\end{block}
		\begin{enumerate}[(A)]
			\item Faster
			\item Backups often contain sensitive data subject to residency laws
			\item Cheaper
			\item Only policy
		\end{enumerate}
	\end{frame}
	
	\begin{frame}{Q337: Explanation}
		\begin{block}{Correct Answer: B}
			\textbf{Backups are data:}
			\begin{itemize}
				\item Sensitive data in backups
				\item Residency applies
				\item Contracts must align
			\end{itemize}
		\end{block}
		\begin{alertblock}{Critical Trap}
			Backups must respect data residency laws.
		\end{alertblock}
	\end{frame}
	
	\begin{frame}{Q338: Data Masking in Testing}
		\begin{block}{Question}
			Why is masking important for non-production testing?
		\end{block}
		\begin{enumerate}[(A)]
			\item Speed only
			\item Protects sensitive production data used in test environments
			\item Cheaper
			\item Less storage
		\end{enumerate}
	\end{frame}
	
	\begin{frame}{Q338: Explanation}
		\begin{block}{Correct Answer: B}
			\textbf{Masking protects sensitive data:}
			\begin{itemize}
				\item Production data in tests
				\item Lower controls in tests
				\item Mask PII and PHI
			\end{itemize}
		\end{block}
		\begin{alertblock}{Critical Trap}
			Test environments often leak unmasked production data.
		\end{alertblock}
	\end{frame}
	
	\begin{frame}{Q339: Data Steward Responsibilities}
		\begin{block}{Question}
			Which is a data steward responsibility?
		\end{block}
		\begin{enumerate}[(A)]
			\item Firewall management
			\item Data quality, metadata, and adherence to data policy
			\item HR
			\item Marketing
		\end{enumerate}
	\end{frame}
	
	\begin{frame}{Q339: Explanation}
		\begin{block}{Correct Answer: B}
			\textbf{Stewardship = quality and policy:}
			\begin{itemize}
				\item Quality
				\item Metadata
				\item Policy adherence
			\end{itemize}
		\end{block}
		\begin{alertblock}{Critical Trap}
			Stewards focus on data, not systems.
		\end{alertblock}
	\end{frame}
	
	\begin{frame}{Q340: Data Retention Metrics}
		\begin{block}{Question}
			Which metric BEST reflects retention compliance?
		\end{block}
		\begin{enumerate}[(A)]
			\item Number of meetings
			\item Percentage of data destroyed per schedule
			\item Marketing spend
			\item Vendor count
		\end{enumerate}
	\end{frame}
	
	\begin{frame}{Q340: Explanation}
		\begin{block}{Correct Answer: B}
			\textbf{Destruction per schedule:}
			\begin{itemize}
				\item Destruction performed per schedule
				\item Exception review
				\item Proof of compliance
			\end{itemize}
		\end{block}
		\begin{alertblock}{Critical Trap}
			Over-retention is a compliance risk.
		\end{alertblock}
	\end{frame}
	
	% ============================================================
	\section{5. SDLC and Emerging Technologies}
	% ============================================================
	
	\begin{frame}{Q341: DevSecOps Culture}
		\begin{block}{Question}
			What cultural change is required for DevSecOps?
		\end{block}
		\begin{enumerate}[(A)]
			\item Security only after deploy
			\item Shared responsibility for security across dev, security, and ops
			\item Fewer tools
			\item More meetings
		\end{enumerate}
	\end{frame}
	
	\begin{frame}{Q341: Explanation}
		\begin{block}{Correct Answer: B}
			\textbf{Shared responsibility:}
			\begin{itemize}
				\item Shared responsibility
				\item Automation
				\item Early feedback
			\end{itemize}
		\end{block}
		\begin{alertblock}{Critical Trap}
			DevSecOps is a culture shift, not a tool.
		\end{alertblock}
	\end{frame}
	
	\begin{frame}{Q342: Supply Chain Software Risk}
		\begin{block}{Question}
			Which is a growing software supply chain risk?
		\end{block}
		\begin{enumerate}[(A)]
			\item Marketing
			\item Malicious or vulnerable third-party dependencies
			\item Larger hard drives
			\item Faster networks
		\end{enumerate}
	\end{frame}
	
	\begin{frame}{Q342: Explanation}
		\begin{block}{Correct Answer: B}
			\textbf{Dependency risk:}
			\begin{itemize}
				\item Open-source dependencies
				\item Malicious packages
				\item SBOM helps
			\end{itemize}
		\end{block}
		\begin{alertblock}{Critical Trap}
			Supply chain software risk is growing.
		\end{alertblock}
	\end{frame}
	
	\begin{frame}{Q343: SBOM Purpose}
		\begin{block}{Question}
			What is an SBOM?
		\end{block}
		\begin{enumerate}[(A)]
			\item Marketing report
			\item Software Bill of Materials listing components
			\item Budget
			\item Vendor logo
		\end{enumerate}
	\end{frame}
	
	\begin{frame}{Q343: Explanation}
		\begin{block}{Correct Answer: B}
			\textbf{SBOM lists components:}
			\begin{itemize}
				\item Components and versions
				\item Supports incident response
				\item Increasingly required by customers
			\end{itemize}
		\end{block}
		\begin{alertblock}{Critical Trap}
			SBOMs speed vulnerability response.
		\end{alertblock}
	\end{frame}
	
	\begin{frame}{Q344: IaC Scanning}
		\begin{block}{Question}
			Why scan Infrastructure as Code?
		\end{block}
		\begin{enumerate}[(A)]
			\item Aesthetics
			\item Detect misconfigurations and embedded secrets before deployment
			\item Faster
			\item Bigger files
		\end{enumerate}
	\end{frame}
	
	\begin{frame}{Q344: Explanation}
		\begin{block}{Correct Answer: B}
			\textbf{Catch issues early:}
			\begin{itemize}
				\item Misconfigurations
				\item Secrets in code
				\item Prevent drift
			\end{itemize}
		\end{block}
		\begin{alertblock}{Critical Trap}
			IaC scanning prevents scaled misconfigurations.
		\end{alertblock}
	\end{frame}
	
	\begin{frame}{Q345: Container Runtime Security}
		\begin{block}{Question}
			Why is container runtime monitoring needed if images are scanned?
		\end{block}
		\begin{enumerate}[(A)]
			\item Cheaper
			\item Detects threats that appear after deployment
			\item Faster
			\item Slower
		\end{enumerate}
	\end{frame}
	
	\begin{frame}{Q345: Explanation}
		\begin{block}{Correct Answer: B}
			\textbf{Image scanning is snapshot:}
			\begin{itemize}
				\item Image scanning is point-in-time
				\item Runtime threats appear later
				\item Monitoring closes gap
			\end{itemize}
		\end{block}
		\begin{alertblock}{Critical Trap}
			Scan at build, monitor at runtime.
		\end{alertblock}
	\end{frame}
	
	\begin{frame}{Q346: API Authentication Strength}
		\begin{block}{Question}
			Which is the STRONGEST API authentication approach?
		\end{block}
		\begin{enumerate}[(A)]
			\item Basic auth in URL
			\item OAuth 2.0 with short-lived tokens and audience restrictions
			\item Shared secret in code
			\item No auth
		\end{enumerate}
	\end{frame}
	
	\begin{frame}{Q346: Explanation}
		\begin{block}{Correct Answer: B}
			\textbf{OAuth 2.0 with short tokens:}
			\begin{itemize}
				\item Short-lived tokens
				\item Audience restriction
				\item Least privilege
			\end{itemize}
		\end{block}
		\begin{alertblock}{Critical Trap}
			Basic auth is weakest; use OAuth 2.0.
		\end{alertblock}
	\end{frame}
	
	\begin{frame}{Q347: Mobile Rooting Risk}
		\begin{block}{Question}
			Why does device rooting or jailbreaking increase risk?
		\end{block}
		\begin{enumerate}[(A)]
			\item Bigger screen
			\item Removes OS security controls and enables unauthorized access
			\item More apps
			\item Bigger battery
		\end{enumerate}
	\end{frame}
	
	\begin{frame}{Q347: Explanation}
		\begin{block}{Correct Answer: B}
			\textbf{Rooting removes controls:}
			\begin{itemize}
				\item Removes OS security controls
				\item Bypasses sandboxing
				\item Increases malware risk
			\end{itemize}
		\end{block}
		\begin{alertblock}{Critical Trap}
			Rooted devices are riskier.
		\end{alertblock}
	\end{frame}
	
	\begin{frame}{Q348: Zero-Trust Policy Engine}
		\begin{block}{Question}
			What role does a policy engine play in Zero Trust?
		\end{block}
		\begin{enumerate}[(A)]
			\item Marketing
			\item Evaluates context and enforces access decisions
			\item Stores backups
			\item Filters email
		\end{enumerate}
	\end{frame}
	
	\begin{frame}{Q348: Explanation}
		\begin{block}{Correct Answer: B}
			\textbf{Policy engine enforces access:}
			\begin{itemize}
				\item Context evaluation
				\item Access decisions
				\item Enforced dynamically
			\end{itemize}
		\end{block}
		\begin{alertblock}{Critical Trap}
			Zero Trust needs a policy engine for consistent decisions.
		\end{alertblock}
	\end{frame}
	
	\begin{frame}{Q349: Post-Quantum Readiness Steps}
		\begin{block}{Question}
			What is the FIRST step in post-quantum readiness?
		\end{block}
		\begin{enumerate}[(A)]
			\item Replace all systems
			\item Inventory cryptographic assets and dependencies
			\item Ignore
			\item Buy tools
		\end{enumerate}
	\end{frame}
	
	\begin{frame}{Q349: Explanation}
		\begin{block}{Correct Answer: B}
			\textbf{Inventory first:}
			\begin{itemize}
				\item Crypto inventory
				\item Dependencies
				\item Migration planning
			\end{itemize}
		\end{block}
		\begin{alertblock}{Critical Trap}
			Quantum readiness begins with inventory.
		\end{alertblock}
	\end{frame}
	
	\begin{frame}{Q350: AI Model Monitoring}
		\begin{block}{Question}
			Why must AI models be monitored after deployment?
		\end{block}
		\begin{enumerate}[(A)]
			\item Cost only
			\item Data drift and behavior changes can reduce accuracy and introduce bias
			\item To satisfy marketing
			\item To delete models
		\end{enumerate}
	\end{frame}
	
	\begin{frame}{Q350: Explanation}
		\begin{block}{Correct Answer: B}
			\textbf{Models drift:}
			\begin{itemize}
				\item Data drift
				\item Behavioral changes
				\item Bias emergence
			\end{itemize}
		\end{block}
		\begin{alertblock}{Critical Trap}
			Model monitoring is essential after deployment.
		\end{alertblock}
	\end{frame}
	
	% ============================================================
	\section{6. Information Security Principles}
	% ============================================================
	
	\begin{frame}{Q351: Non-Repudiation Evidence}
		\begin{block}{Question}
			Which control provides the STRONGEST non-repudiation?
		\end{block}
		\begin{enumerate}[(A)]
			\item Password
			\item Digital signature with immutable logs
			\item Session cookie
			\item Firewall rule
		\end{enumerate}
	\end{frame}
	
	\begin{frame}{Q351: Explanation}
		\begin{block}{Correct Answer: B}
			\textbf{Signatures and logs support proof:}
			\begin{itemize}
				\item Signed with private key
				\item Logged immutably
				\item Provides proof
			\end{itemize}
		\end{block}
		\begin{alertblock}{Critical Trap}
			Non-repudiation requires cryptographic proof.
		\end{alertblock}
	\end{frame}
	
	\begin{frame}{Q352: TLS Termination Risk}
		\begin{block}{Question}
			What is the risk of terminating TLS at a load balancer without re-encryption?
		\end{block}
		\begin{enumerate}[(A)]
			\item Slower
			\item Traffic between load balancer and backend may be exposed
			\item Cheaper
			\item Faster
		\end{enumerate}
	\end{frame}
	
	\begin{frame}{Q352: Explanation}
		\begin{block}{Correct Answer: B}
			\textbf{Re-encrypt or segment:}
			\begin{itemize}
				\item Exposed backend
				\item Re-encrypt or segment
				\item Verify with threat model
			\end{itemize}
		\end{block}
		\begin{alertblock}{Critical Trap}
			TLS termination exposes backend if not re-encrypted.
		\end{alertblock}
	\end{frame}
	
	\begin{frame}{Q353: Certificate Transparency}
		\begin{block}{Question}
			What does certificate transparency provide?
		\end{block}
		\begin{enumerate}[(A)]
			\item Marketing
			\item Public logs enabling detection of misissued certificates
			\item Bigger certificates
			\item Cheaper TLS
		\end{enumerate}
	\end{frame}
	
	\begin{frame}{Q353: Explanation}
		\begin{block}{Correct Answer: B}
			\textbf{Public logs detect misissuance:}
			\begin{itemize}
				\item Public logs
				\item Detect misissuance
				\item Early warning
			\end{itemize}
		\end{block}
		\begin{alertblock}{Critical Trap}
			Certificate Transparency helps detect fraud.
		\end{alertblock}
	\end{frame}
	
	\begin{frame}{Q354: HSM Tamper Resistance}
		\begin{block}{Question}
			Why does an HSM provide stronger protection for keys than a server?
		\end{block}
		\begin{enumerate}[(A)]
			\item Bigger keys
			\item Tamper-resistant hardware, keys never leave in plaintext
			\item Cheaper
			\item Faster
		\end{enumerate}
	\end{frame}
	
	\begin{frame}{Q354: Explanation}
		\begin{block}{Correct Answer: B}
			\textbf{Tamper resistance:}
			\begin{itemize}
				\item Tamper-resistant
				\item Keys do not leave
				\item Operations inside
			\end{itemize}
		\end{block}
		\begin{alertblock}{Critical Trap}
			HSMs protect keys even if host is compromised.
		\end{alertblock}
	\end{frame}
	
	\begin{frame}{Q355: Symmetric vs Asymmetric Use}
		\begin{block}{Question}
			Which is a correct statement?
		\end{block}
		\begin{enumerate}[(A)]
			\item Symmetric is used for key exchange only
			\item Asymmetric is used for key exchange; symmetric for bulk encryption
			\item Asymmetric for bulk
			\item Neither
		\end{enumerate}
	\end{frame}
	
	\begin{frame}{Q355: Explanation}
		\begin{block}{Correct Answer: B}
			\textbf{Common hybrid scheme:}
			\begin{itemize}
				\item Asymmetric for key exchange
				\item Symmetric for bulk
				\item Common hybrid scheme
			\end{itemize}
		\end{block}
		\begin{alertblock}{Critical Trap}
			Asymmetric establishes; symmetric encrypts bulk.
		\end{alertblock}
	\end{frame}
	
	\begin{frame}{Q356: Digital Signature Verification}
		\begin{block}{Question}
			Which key is used to verify a digital signature?
		\end{block}
		\begin{enumerate}[(A)]
			\item Sender private key
			\item Sender public key
			\item Receiver private key
			\item Receiver public key
		\end{enumerate}
	\end{frame}
	
	\begin{frame}{Q356: Explanation}
		\begin{block}{Correct Answer: B}
			\textbf{Sender public key:}
			\begin{itemize}
				\item Verified with sender public key
				\item Signed with sender private key
				\item Proves origin
			\end{itemize}
		\end{block}
		\begin{alertblock}{Critical Trap}
			Verify with sender public key.
		\end{alertblock}
	\end{frame}
	
	\begin{frame}{Q357: Access Review Metrics}
		\begin{block}{Question}
			Which metric shows access review effectiveness?
		\end{block}
		\begin{enumerate}[(A)]
			\item Number of meetings
			\item Percentage of access reviewed and revoked where not needed
			\item Vendor count
			\item Budget
		\end{enumerate}
	\end{frame}
	
	\begin{frame}{Q357: Explanation}
		\begin{block}{Correct Answer: B}
			\textbf{Review completion plus revocation:}
			\begin{itemize}
				\item Review completion
				\item Revocation
				\item Reduced access creep
			\end{itemize}
		\end{block}
		\begin{alertblock}{Critical Trap}
			Reviews without revocations do little.
		\end{alertblock}
	\end{frame}
	
	\begin{frame}{Q358: UEBA False Positive Risk}
		\begin{block}{Question}
			What is the main risk of poorly tuned UEBA?
		\end{block}
		\begin{enumerate}[(A)]
			\item Too few alerts
			\item Alert fatigue from false positives
			\item Slower system
			\item Smaller logs
		\end{enumerate}
	\end{frame}
	
	\begin{frame}{Q358: Explanation}
		\begin{block}{Correct Answer: B}
			\textbf{Tune to reduce noise:}
			\begin{itemize}
				\item False positives
				\item Alert fatigue
				\item Tuning essential
			\end{itemize}
		\end{block}
		\begin{alertblock}{Critical Trap}
			Untuned UEBA causes alert fatigue.
		\end{alertblock}
	\end{frame}
	
	\begin{frame}{Q359: Federation Trust Risk}
		\begin{block}{Question}
			What is the primary risk in identity federation?
		\end{block}
		\begin{enumerate}[(A)]
			\item Slower logins
			\item Compromise of a federation partner can grant access to your systems
			\item Cheaper
			\item More tools
		\end{enumerate}
	\end{frame}
	
	\begin{frame}{Q359: Explanation}
		\begin{block}{Correct Answer: B}
			\textbf{Federation partner trust:}
			\begin{itemize}
				\item Partner compromise
				\item Trust propagation
				\item Contracts and monitoring
			\end{itemize}
		\end{block}
		\begin{alertblock}{Critical Trap}
			Federated trust amplifies risk from partners.
		\end{alertblock}
	\end{frame}
	
	\begin{frame}{Q360: Security Metrics Balance}
		\begin{block}{Question}
			Which metric combination BEST reflects security program health?
		\end{block}
		\begin{enumerate}[(A)]
			\item Only number of incidents
			\item Mix of risk reduction, control effectiveness, and outcome metrics
			\item Only budget
			\item Only training completion
		\end{enumerate}
	\end{frame}
	
	\begin{frame}{Q360: Explanation}
		\begin{block}{Correct Answer: B}
			\textbf{Balanced metrics:}
			\begin{itemize}
				\item Risk reduction
				\item Control effectiveness
				\item Outcomes
			\end{itemize}
		\end{block}
		\begin{alertblock}{Critical Trap}
			Balanced metrics tell the real story.
		\end{alertblock}
	\end{frame}
	
	\begin{frame}{Q361: Alert Prioritization}
		\begin{block}{Question}
			How should security alerts be prioritized?
		\end{block}
		\begin{enumerate}[(A)]
			\item Alphabetically
			\item By business impact and asset criticality
			\item By age
			\item By tool
		\end{enumerate}
	\end{frame}
	
	\begin{frame}{Q361: Explanation}
		\begin{block}{Correct Answer: B}
			\textbf{Business impact drives priority:}
			\begin{itemize}
				\item Business impact
				\item Asset criticality
				\item Alert context
			\end{itemize}
		\end{block}
		\begin{alertblock}{Critical Trap}
			Prioritize by impact, not by ticket age.
		\end{alertblock}
	\end{frame}
	
	\begin{frame}{Q362: Threat Hunting Hypothesis}
		\begin{block}{Question}
			What starts a threat hunt?
		\end{block}
		\begin{enumerate}[(A)]
			\item Marketing
			\item A hypothesis based on threat intel or anomalies
			\item Random scanning
			\item Vendor discount
		\end{enumerate}
	\end{frame}
	
	\begin{frame}{Q362: Explanation}
		\begin{block}{Correct Answer: B}
			\textbf{Hypothesis-driven:}
			\begin{itemize}
				\item Hypothesis
				\item Intel-driven
				\item Anomalies
			\end{itemize}
		\end{block}
		\begin{alertblock}{Critical Trap}
			Hunting is not random; it is targeted.
		\end{alertblock}
	\end{frame}
	
	\begin{frame}{Q363: Threat Intel Validation}
		\begin{block}{Question}
			Why should threat intelligence indicators be validated?
		\end{block}
		\begin{enumerate}[(A)]
			\item Marketing
			\item False positives and stale data cause unnecessary blocks and alert fatigue
			\item Cheaper
			\item Faster
		\end{enumerate}
	\end{frame}
	
	\begin{frame}{Q363: Explanation}
		\begin{block}{Correct Answer: B}
			\textbf{Validate before blocking:}
			\begin{itemize}
				\item False positives
				\item Stale indicators
				\item Business disruption risk
			\end{itemize}
		\end{block}
		\begin{alertblock}{Critical Trap}
			Unvalidated intel harms business.
		\end{alertblock}
	\end{frame}
	
	\begin{frame}{Q364: IR Tabletop Content}
		\begin{block}{Question}
			What should an IR tabletop exercise cover?
		\end{block}
		\begin{enumerate}[(A)]
			\item Only technology
			\item Decision-making, communications, roles, and escalation
			\item Only backups
			\item Only vendor contacts
		\end{enumerate}
	\end{frame}
	
	\begin{frame}{Q364: Explanation}
		\begin{block}{Correct Answer: B}
			\textbf{Cover decisions and comms:}
			\begin{itemize}
				\item Decisions
				\item Communications
				\item Roles
				\item Escalation
			\end{itemize}
		\end{block}
		\begin{alertblock}{Critical Trap}
			Tabletop tests decisions, not just technology.
		\end{alertblock}
	\end{frame}
	
	\begin{frame}{Q365: Post-Incident Accountability}
		\begin{block}{Question}
			What should accompany post-incident action items?
		\end{block}
		\begin{enumerate}[(A)]
			\item Blame
			\item Named owners and due dates
			\item Budget increase
			\item Marketing
		\end{enumerate}
	\end{frame}
	
	\begin{frame}{Q365: Explanation}
		\begin{block}{Correct Answer: B}
			\textbf{Owners and dates:}
			\begin{itemize}
				\item Owners
				\item Due dates
				\item Tracking
			\end{itemize}
		\end{block}
		\begin{alertblock}{Critical Trap}
			Actions without owners are wishes.
		\end{alertblock}
	\end{frame}
	
	\begin{frame}{Q366: Forensic Chain of Custody}
		\begin{block}{Question}
			What does chain of custody provide?
		\end{block}
		\begin{enumerate}[(A)]
			\item Marketing
			\item Documented evidence integrity from collection to presentation
			\item Cheaper
			\item Faster
		\end{enumerate}
	\end{frame}
	
	\begin{frame}{Q366: Explanation}
		\begin{block}{Correct Answer: B}
			\textbf{Custody preserves admissibility:}
			\begin{itemize}
				\item Documented integrity
				\item Admissibility
				\item Regulatory support
			\end{itemize}
		\end{block}
		\begin{alertblock}{Critical Trap}
			Broken chain of custody loses evidence.
		\end{alertblock}
	\end{frame}
	
	\begin{frame}{Q367: Vendor Breach Response}
		\begin{block}{Question}
			What should be done IMMEDIATELY after a vendor notifies of a breach?
		\end{block}
		\begin{enumerate}[(A)]
			\item Terminate the vendor
			\item Verify scope, impact, and applicable notification requirements
			\item Blame
			\item Ignore
		\end{enumerate}
	\end{frame}
	
	\begin{frame}{Q367: Explanation}
		\begin{block}{Correct Answer: B}
			\textbf{Verify and respond:}
			\begin{itemize}
				\item Verify scope
				\item Assess impact
				\item Regulatory notification
			\end{itemize}
		\end{block}
		\begin{alertblock}{Critical Trap}
			Vendor breach is your breach to manage.
		\end{alertblock}
	\end{frame}
	
	\begin{frame}{Q368: Patch Prioritization}
		\begin{block}{Question}
			Which patch should be prioritized?
		\end{block}
		\begin{enumerate}[(A)]
			\item Oldest patch
			\item Known exploited vulnerability on an internet-facing critical system
			\item Newest patch
			\item Vendor preferred
		\end{enumerate}
	\end{frame}
	
	\begin{frame}{Q368: Explanation}
		\begin{block}{Correct Answer: B}
			\textbf{Exploited critical:}
			\begin{itemize}
				\item Known exploited
				\item Internet-facing
				\item Critical asset
			\end{itemize}
		\end{block}
		\begin{alertblock}{Critical Trap}
			Priority = exploitability times impact.
		\end{alertblock}
	\end{frame}
	
	\begin{frame}{Q369: Secure Baseline Drift}
		\begin{block}{Question}
			What is configuration drift and why does it matter?
		\end{block}
		\begin{enumerate}[(A)]
			\item Marketing
			\item Systems diverge from secure baseline, creating exposure
			\item Faster
			\item Smaller
		\end{enumerate}
	\end{frame}
	
	\begin{frame}{Q369: Explanation}
		\begin{block}{Correct Answer: B}
			\textbf{Drift creates exposure:}
			\begin{itemize}
				\item Divergence
				\item Detection
				\item Remediation
			\end{itemize}
		\end{block}
		\begin{alertblock}{Critical Trap}
			Configuration drift is silent risk.
		\end{alertblock}
	\end{frame}
	
	\begin{frame}{Q370: Backup vs Archive Retention}
		\begin{block}{Question}
			Which retention policy applies to backups versus archives?
		\end{block}
		\begin{enumerate}[(A)]
			\item Same
			\item Backup retention is short-term; archive retention follows regulatory requirements
			\item Archive short-term
			\item Neither
		\end{enumerate}
	\end{frame}
	
	\begin{frame}{Q370: Explanation}
		\begin{block}{Correct Answer: B}
			\textbf{Different retention:}
			\begin{itemize}
				\item Backup retention short
				\item Archive follows regulation
				\item Distinguish policies
			\end{itemize}
		\end{block}
		\begin{alertblock}{Critical Trap}
			Backups and archives have different retention.
		\end{alertblock}
	\end{frame}
	
	\begin{frame}{Q371: Immutable Logs Risk}
		\begin{block}{Question}
			What is a trade-off of immutable logs?
		\end{block}
		\begin{enumerate}[(A)]
			\item Slower search
			\item Storage cost growth; sensitive data persists in logs
			\item Cheaper
			\item Faster
		\end{enumerate}
	\end{frame}
	
	\begin{frame}{Q371: Explanation}
		\begin{block}{Correct Answer: B}
			\textbf{Storage and sensitive data:}
			\begin{itemize}
				\item Storage growth
				\item Sensitive data persists
				\item Retention policies
			\end{itemize}
		\end{block}
		\begin{alertblock}{Critical Trap}
			Immutable logs accumulate data; retention management matters.
		\end{alertblock}
	\end{frame}
	
	\begin{frame}{Q372: Key Rotation Frequency}
		\begin{block}{Question}
			How should key rotation frequency be determined?
		\end{block}
		\begin{enumerate}[(A)]
			\item Never rotate
			\item Based on crypto period, data sensitivity, and regulatory requirements
			\item Once per year only
			\item Vendor preference
		\end{enumerate}
	\end{frame}
	
	\begin{frame}{Q372: Explanation}
		\begin{block}{Correct Answer: B}
			\textbf{Risk-based rotation:}
			\begin{itemize}
				\item Crypto period
				\item Sensitivity
				\item Regulation
			\end{itemize}
		\end{block}
		\begin{alertblock}{Critical Trap}
			Rotation frequency is risk and regulator driven.
		\end{alertblock}
	\end{frame}
	
	\begin{frame}{Q373: HSM Availability Risk}
		\begin{block}{Question}
			What is the PRIMARY risk of a single HSM without high availability?
		\end{block}
		\begin{enumerate}[(A)]
			\item Cheaper
			\item HSM failure can halt cryptographic operations and services
			\item Slower
			\item Smaller
		\end{enumerate}
	\end{frame}
	
	\begin{frame}{Q373: Explanation}
		\begin{block}{Correct Answer: B}
			\textbf{HA matters:}
			\begin{itemize}
				\item Crypto becomes unavailable
				\item Service outage
				\item HA design
			\end{itemize}
		\end{block}
		\begin{alertblock}{Critical Trap}
			Single HSM is a single point of failure.
		\end{alertblock}
	\end{frame}
	
	\begin{frame}{Q374: TLS Version}
		\begin{block}{Question}
			Which TLS version should be avoided?
		\end{block}
		\begin{enumerate}[(A)]
			\item TLS 1.3
			\item TLS 1.0 and 1.1
			\item TLS 1.2
			\item All versions
		\end{enumerate}
	\end{frame}
	
	\begin{frame}{Q374: Explanation}
		\begin{block}{Correct Answer: B}
			\textbf{Old TLS deprecated:}
			\begin{itemize}
				\item TLS 1.0 and 1.1 deprecated
				\item Use TLS 1.2 or 1.3
				\item Disable weak versions
			\end{itemize}
		\end{block}
		\begin{alertblock}{Critical Trap}
			TLS 1.0 and 1.1 are insecure.
		\end{alertblock}
	\end{frame}
	
	\begin{frame}{Q375: Certificate Renewal Monitoring}
		\begin{block}{Question}
			Why is certificate renewal monitoring critical?
		\end{block}
		\begin{enumerate}[(A)]
			\item Marketing
			\item Prevent outages from expiring certificates and detect unauthorized changes
			\item Cheaper
			\item Faster
		\end{enumerate}
	\end{frame}
	
	\begin{frame}{Q375: Explanation}
		\begin{block}{Correct Answer: B}
			\textbf{Expiry causes outages:}
			\begin{itemize}
				\item Expiry causes outages
				\item Monitoring alerts
				\item Renewal plan
			\end{itemize}
		\end{block}
		\begin{alertblock}{Critical Trap}
			Expired certificates break services.
		\end{alertblock}
	\end{frame}
	
	\begin{frame}{Q376: HSM vs TPM}
		\begin{block}{Question}
			What is the main difference between HSM and TPM?
		\end{block}
		\begin{enumerate}[(A)]
			\item Identical
			\item TPM is embedded in device for device identity; HSM is dedicated to server or enterprise key management
			\item TPM is for servers only
			\item HSM is for laptops
		\end{enumerate}
	\end{frame}
	
	\begin{frame}{Q376: Explanation}
		\begin{block}{Correct Answer: B}
			\textbf{Device vs enterprise:}
			\begin{itemize}
				\item TPM in endpoint devices
				\item HSM in server or enterprise
				\item Different scale
			\end{itemize}
		\end{block}
		\begin{alertblock}{Critical Trap}
			TPM is per-device; HSM is enterprise-grade.
		\end{alertblock}
	\end{frame}
	
	\begin{frame}{Q377: Access Recertification Ownership}
		\begin{block}{Question}
			Who should approve access recertification?
		\end{block}
		\begin{enumerate}[(A)]
			\item IT
			\item Business managers who know the user's role
			\item Auditor
			\item Vendor
		\end{enumerate}
	\end{frame}
	
	\begin{frame}{Q377: Explanation}
		\begin{block}{Correct Answer: B}
			\textbf{Business managers know roles:}
			\begin{itemize}
				\item Business managers
				\item Knowing role
				\item Accountability
			\end{itemize}
		\end{block}
		\begin{alertblock}{Critical Trap}
			Recertification is a business decision.
		\end{alertblock}
	\end{frame}
	
	\begin{frame}{Q378: Least Privilege Violation}
		\begin{block}{Question}
			What is the PRIMARY risk of granting excessive privileges?
		\end{block}
		\begin{enumerate}[(A)]
			\item Marketing
			\item Blast radius of misuse or compromise
			\item Cheaper
			\item Faster
		\end{enumerate}
	\end{frame}
	
	\begin{frame}{Q378: Explanation}
		\begin{block}{Correct Answer: B}
			\textbf{Blast radius grows:}
			\begin{itemize}
				\item Blast radius grows
				\item Misuse impact
				\item Compromise impact
			\end{itemize}
		\end{block}
		\begin{alertblock}{Critical Trap}
			Excessive privilege magnifies impact.
		\end{alertblock}
	\end{frame}
	
	\begin{frame}{Q379: Privacy by Design Benefit}
		\begin{block}{Question}
			What is a benefit of privacy by design?
		\end{block}
		\begin{enumerate}[(A)]
			\item Higher cost
			\item Cheaper remediation later and improved trust
			\item Slower
			\item More audits
		\end{enumerate}
	\end{frame}
	
	\begin{frame}{Q379: Explanation}
		\begin{block}{Correct Answer: B}
			\textbf{Early design cheaper:}
			\begin{itemize}
				\item Cheaper later
				\item Improved trust
				\item Compliance built in
			\end{itemize}
		\end{block}
		\begin{alertblock}{Critical Trap}
			Privacy by design saves cost and builds trust.
		\end{alertblock}
	\end{frame}
	
	\begin{frame}{Q380: DSAR Refusal Documentation}
		\begin{block}{Question}
			What should be documented when refusing a data subject request?
		\end{block}
		\begin{enumerate}[(A)]
			\item Nothing
			\item The legal basis and communication to the data subject
			\item Marketing
			\item Budget
		\end{enumerate}
	\end{frame}
	
	\begin{frame}{Q380: Explanation}
		\begin{block}{Correct Answer: B}
			\textbf{Document refusal basis:}
			\begin{itemize}
				\item Legal basis
				\item Clear communication
				\item Records
			\end{itemize}
		\end{block}
		\begin{alertblock}{Critical Trap}
			Refusals need documented basis.
		\end{alertblock}
	\end{frame}
	
	\begin{frame}{Q381: HIPAA Safeguards}
		\begin{block}{Question}
			Which are HIPAA safeguard categories?
		\end{block}
		\begin{enumerate}[(A)]
			\item Marketing
			\item Administrative, physical, and technical safeguards
			\item Only technical
			\item Only physical
		\end{enumerate}
	\end{frame}
	
	\begin{frame}{Q381: Explanation}
		\begin{block}{Correct Answer: B}
			\textbf{Three safeguard types:}
			\begin{itemize}
				\item Administrative
				\item Physical
				\item Technical
			\end{itemize}
		\end{block}
		\begin{alertblock}{Critical Trap}
			HIPAA security requires all three safeguard types.
		\end{alertblock}
	\end{frame}
	
	\begin{frame}{Q382: CCPA Verification}
		\begin{block}{Question}
			What is a key requirement for CCPA consumer requests?
		\end{block}
		\begin{enumerate}[(A)]
			\item Instant response
			\item Verify the consumer and provide a reasonable response process
			\item Ignore
			\item Marketing
		\end{enumerate}
	\end{frame}
	
	\begin{frame}{Q382: Explanation}
		\begin{block}{Correct Answer: B}
			\textbf{Verify and respond:}
			\begin{itemize}
				\item Verify consumer
				\item Reasonable process
				\item Documented handling
			\end{itemize}
		\end{block}
		\begin{alertblock}{Critical Trap}
			CCPA requests require verification.
		\end{alertblock}
	\end{frame}
	
	\begin{frame}{Q383: Encryption Algorithm Choice}
		\begin{block}{Question}
			Which encryption algorithm is considered strong today?
		\end{block}
		\begin{enumerate}[(A)]
			\item DES
			\item AES-256
			\item RC4
			\item MD5
		\end{enumerate}
	\end{frame}
	
	\begin{frame}{Q383: Explanation}
		\begin{block}{Correct Answer: B}
			\textbf{AES-256 strong:}
			\begin{itemize}
				\item AES-256 recommended
				\item DES, RC4 deprecated
				\item MD5 is a hash
			\end{itemize}
		\end{block}
		\begin{alertblock}{Critical Trap}
			DES and RC4 are outdated.
		\end{alertblock}
	\end{frame}
	
	\begin{frame}{Q384: Signature vs MAC}
		\begin{block}{Question}
			What is the main difference between a digital signature and a MAC?
		\end{block}
		\begin{enumerate}[(A)]
			\item Identical
			\item Signature uses asymmetric keys and provides non-repudiation; MAC uses shared key only
			\item MAC uses asymmetric
			\item Both provide non-repudiation
		\end{enumerate}
	\end{frame}
	
	\begin{frame}{Q384: Explanation}
		\begin{block}{Correct Answer: B}
			\textbf{Signature vs MAC:}
			\begin{itemize}
				\item Signature asymmetric
				\item MAC symmetric
				\item Only signature gives non-repudiation
			\end{itemize}
		\end{block}
		\begin{alertblock}{Critical Trap}
			MAC cannot provide non-repudiation.
		\end{alertblock}
	\end{frame}
	
	\begin{frame}{Q385: Security Training Metrics}
		\begin{block}{Question}
			Which training metric is MOST useful?
		\end{block}
		\begin{enumerate}[(A)]
			\item Completion rate
			\item Behavior change such as phishing reports and reduced clicks
			\item Number of slides
			\item Hours trained
		\end{enumerate}
	\end{frame}
	
	\begin{frame}{Q385: Explanation}
		\begin{block}{Correct Answer: B}
			\textbf{Behavior beats completion:}
			\begin{itemize}
				\item Behavior metrics
				\item Phishing reports
				\item Reduced clicks
			\end{itemize}
		\end{block}
		\begin{alertblock}{Critical Trap}
			Completion alone does not equal awareness.
		\end{alertblock}
	\end{frame}
	
	\begin{frame}{Q386: Whaling Indicators}
		\begin{block}{Question}
			Which is a hallmark of a whaling attack?
		\end{block}
		\begin{enumerate}[(A)]
			\item Mass phishing
			\item Highly personalized messages targeting executives
			\item Marketing
			\item Only spam
		\end{enumerate}
	\end{frame}
	
	\begin{frame}{Q386: Explanation}
		\begin{block}{Correct Answer: B}
			\textbf{Highly targeted:}
			\begin{itemize}
				\item Personalized
				\item Executives targeted
				\item Spear phishing variant
			\end{itemize}
		\end{block}
		\begin{alertblock}{Critical Trap}
			Whaling is highly targeted.
		\end{alertblock}
	\end{frame}
	
	\begin{frame}{Q387: Board Tabletop Outcome}
		\begin{block}{Question}
			Which outcome is MOST valuable from a board cyber tabletop?
		\end{block}
		\begin{enumerate}[(A)]
			\item Slide count
			\item Clarity of decision rights and crisis communication expectations
			\item Marketing
			\item Budget
		\end{enumerate}
	\end{frame}
	
	\begin{frame}{Q387: Explanation}
		\begin{block}{Correct Answer: B}
			\textbf{Clarity of roles:}
			\begin{itemize}
				\item Decision rights
				\item Communication clarity
				\item Board roles
			\end{itemize}
		\end{block}
		\begin{alertblock}{Critical Trap}
			Board tabletops clarify crisis roles.
		\end{alertblock}
	\end{frame}
	
	\begin{frame}{Q388: Middleware Patch Risk}
		\begin{block}{Question}
			Why is unpatched middleware high-risk?
		\end{block}
		\begin{enumerate}[(A)]
			\item Cheap
			\item Sits between tiers and exposes high-value systems if compromised
			\item Slower
			\item Bigger
		\end{enumerate}
	\end{frame}
	
	\begin{frame}{Q388: Explanation}
		\begin{block}{Correct Answer: B}
			\textbf{High-value position:}
			\begin{itemize}
				\item Middle tier
				\item Broad impact
				\item High-value position
			\end{itemize}
		\end{block}
		\begin{alertblock}{Critical Trap}
			Middleware vulnerabilities have amplified impact.
		\end{alertblock}
	\end{frame}
	
	\begin{frame}{Q389: Secret Storage Practice}
		\begin{block}{Question}
			Which is the best practice for secret storage?
		\end{block}
		\begin{enumerate}[(A)]
			\item Hardcode
			\item Secrets manager with rotation and audit
			\item Email
			\item Shared document
		\end{enumerate}
	\end{frame}
	
	\begin{frame}{Q389: Explanation}
		\begin{block}{Correct Answer: B}
			\textbf{Secrets manager:}
			\begin{itemize}
				\item Central storage
				\item Rotation
				\item Audit
			\end{itemize}
		\end{block}
		\begin{alertblock}{Critical Trap}
			Hardcoded secrets are a breach risk.
		\end{alertblock}
	\end{frame}
	
	\begin{frame}{Q390: Continuous Compliance Value}
		\begin{block}{Question}
			What is the main value of continuous compliance?
		\end{block}
		\begin{enumerate}[(A)]
			\item Marketing
			\item Detects drift and non-compliance faster than periodic audits
			\item Cheaper
			\item Faster
		\end{enumerate}
	\end{frame}
	
	\begin{frame}{Q390: Explanation}
		\begin{block}{Correct Answer: B}
			\textbf{Continuous detection:}
			\begin{itemize}
				\item Faster detection
				\item Drift remediation
				\item Continuous assurance
			\end{itemize}
		\end{block}
		\begin{alertblock}{Critical Trap}
			Continuous compliance closes audit gaps.
		\end{alertblock}
	\end{frame}
	
	\begin{frame}{Q391: CAIQ Limitations}
		\begin{block}{Question}
			What is a limitation of CAIQ?
		\end{block}
		\begin{enumerate}[(A)]
			\item Independent audit
			\item Self-assessment may be biased or inaccurate
			\item Legal evidence
			\item Regulatory filing
		\end{enumerate}
	\end{frame}
	
	\begin{frame}{Q391: Explanation}
		\begin{block}{Correct Answer: B}
			\textbf{Self-reported:}
			\begin{itemize}
				\item Self-reported
				\item Bias risk
				\item Independent audits stronger
			\end{itemize}
		\end{block}
		\begin{alertblock}{Critical Trap}
			CAIQ is not independent assurance.
		\end{alertblock}
	\end{frame}
	
	\begin{frame}{Q392: SaaS Misconfiguration Prevention}
		\begin{block}{Question}
			What reduces SaaS misconfiguration risk?
		\end{block}
		\begin{enumerate}[(A)]
			\item Nothing
			\item Continuous configuration monitoring and least privilege defaults
			\item Vendor branding
			\item More users
		\end{enumerate}
	\end{frame}
	
	\begin{frame}{Q392: Explanation}
		\begin{block}{Correct Answer: B}
			\textbf{Monitor configurations:}
			\begin{itemize}
				\item Configuration monitoring
				\item Least privilege
				\item User training
			\end{itemize}
		\end{block}
		\begin{alertblock}{Critical Trap}
			Misconfigurations are the top SaaS risk.
		\end{alertblock}
	\end{frame}
	
	\begin{frame}{Q393: Cloud IAM MFA}
		\begin{block}{Question}
			Why is MFA essential for cloud IAM?
		\end{block}
		\begin{enumerate}[(A)]
			\item Marketing
			\item Prevents credential compromise and unauthorized access to cloud resources
			\item Cheaper
			\item Faster
		\end{enumerate}
	\end{frame}
	
	\begin{frame}{Q393: Explanation}
		\begin{block}{Correct Answer: B}
			\textbf{MFA blocks credential abuse:}
			\begin{itemize}
				\item Prevents credential theft
				\item Requires second factor
				\item Reduces cloud account takeover
			\end{itemize}
		\end{block}
		\begin{alertblock}{Critical Trap}
			Cloud without MFA is highly risky.
		\end{alertblock}
	\end{frame}
	
	\begin{frame}{Q394: Cloud Encryption Key Control}
		\begin{block}{Question}
			What does BYOK provide in cloud encryption?
		\end{block}
		\begin{enumerate}[(A)]
			\item Marketing
			\item Customer-controlled key material with provider-managed operations
			\item Free tier
			\item Unlimited storage
		\end{enumerate}
	\end{frame}
	
	\begin{frame}{Q394: Explanation}
		\begin{block}{Correct Answer: B}
			\textbf{BYOK keeps key control:}
			\begin{itemize}
				\item Bring your own key
				\item Control key material
				\item Provider operations
			\end{itemize}
		\end{block}
		\begin{alertblock}{Critical Trap}
			BYOK improves key control.
		\end{alertblock}
	\end{frame}
	
	\begin{frame}{Q395: Cloud Data Loss Prevention}
		\begin{block}{Question}
			Which control helps prevent accidental cloud data deletion?
		\end{block}
		\begin{enumerate}[(A)]
			\item Marketing
			\item Versioning, soft delete, and least privilege
			\item Cheaper
			\item Faster
		\end{enumerate}
	\end{frame}
	
	\begin{frame}{Q395: Explanation}
		\begin{block}{Correct Answer: B}
			\textbf{Versioning and soft delete:}
			\begin{itemize}
				\item Versioning
				\item Soft delete
				\item Least privilege
			\end{itemize}
		\end{block}
		\begin{alertblock}{Critical Trap}
			Cloud data loss is often customer-caused.
		\end{alertblock}
	\end{frame}
	
	\begin{frame}{Q396: Security Metrics Context}
		\begin{block}{Question}
			What makes security metrics valuable to executives?
		\end{block}
		\begin{enumerate}[(A)]
			\item Raw counts
			\item Business context and risk appetite comparison
			\item Technical jargon
			\item Marketing
		\end{enumerate}
	\end{frame}
	
	\begin{frame}{Q396: Explanation}
		\begin{block}{Correct Answer: B}
			\textbf{Business context:}
			\begin{itemize}
				\item Business context
				\item Appetite comparison
				\item Actionable insights
			\end{itemize}
		\end{block}
		\begin{alertblock}{Critical Trap}
			Executives need context, not raw data.
		\end{alertblock}
	\end{frame}
	
	\begin{frame}{Q397: Security Risk Acceptance}
		\begin{block}{Question}
			Who should sign off on a residual security risk that exceeds appetite?
		\end{block}
		\begin{enumerate}[(A)]
			\item Security team
			\item Board or delegated authority
			\item Vendor
			\item Auditor
		\end{enumerate}
	\end{frame}
	
	\begin{frame}{Q397: Explanation}
		\begin{block}{Correct Answer: B}
			\textbf{Board-level for high risk:}
			\begin{itemize}
				\item Board or delegated authority
				\item Documented
				\item Time-bound
			\end{itemize}
		\end{block}
		\begin{alertblock}{Critical Trap}
			Above-appetite acceptance requires high authority.
		\end{alertblock}
	\end{frame}
	
	\begin{frame}{Q398: Program Maturity Indicators}
		\begin{block}{Question}
			Which is a key indicator of program maturity?
		\end{block}
		\begin{enumerate}[(A)]
			\item More tools
			\item Risk-based decisions, metrics, and continuous improvement
			\item Larger budget
			\item Fewer incidents
		\end{enumerate}
	\end{frame}
	
	\begin{frame}{Q398: Explanation}
		\begin{block}{Correct Answer: B}
			\textbf{Maturity is outcomes:}
			\begin{itemize}
				\item Risk-based
				\item Measured
				\item Improving
			\end{itemize}
		\end{block}
		\begin{alertblock}{Critical Trap}
			Maturity is continuous improvement.
		\end{alertblock}
	\end{frame}
	
	\begin{frame}{Q399: Security Continuous Cycle}
		\begin{block}{Question}
			Why is security described as continuous?
		\end{block}
		\begin{enumerate}[(A)]
			\item Marketing
			\item Threats and business change require ongoing risk management
			\item One-time project
			\item Purchase
		\end{enumerate}
	\end{frame}
	
	\begin{frame}{Q399: Explanation}
		\begin{block}{Correct Answer: B}
			\textbf{Continuous risk management:}
			\begin{itemize}
				\item Threats change
				\item Business changes
				\item Ongoing
			\end{itemize}
		\end{block}
		\begin{alertblock}{Critical Trap}
			Security is a continuous process, not a project.
		\end{alertblock}
	\end{frame}
	
	\begin{frame}{Q400: Ultimate Security Truth}
		\begin{block}{Question}
			What is the ultimate truth about security risk?
		\end{block}
		\begin{enumerate}[(A)]
			\item Can be eliminated
			\item Risk is managed to acceptable levels and accepted by owners
			\item One-time decision
			\item Optional
		\end{enumerate}
	\end{frame}
	
	\begin{frame}{Q400: Explanation}
		\begin{block}{Correct Answer: B}
			\textbf{Risk managed, not eliminated:}
			\begin{itemize}
				\item Managed to acceptable levels
				\item Accepted by owners
				\item Continuous
			\end{itemize}
		\end{block}
		\begin{alertblock}{Critical Trap}
			Risk cannot be eliminated; it can be managed.
		\end{alertblock}
	\end{frame}
	
	\begin{frame}{Exam Mindset Summary}
		\begin{block}{Remember}
			\begin{itemize}
				\item \textbf{CIA triad:} Confidentiality, Integrity, Availability
				\item \textbf{IAAA:} Identification, Authentication, Authorization, Accountability
				\item \textbf{Encryption is reversible; hashing is one-way}
				\item \textbf{Signatures} provide integrity, authenticity, non-repudiation
				\item \textbf{Privacy} is rights; \textbf{security} is protection
				\item \textbf{Zero Trust:} verify every access
				\item \textbf{Cloud:} shared responsibility; data location matters
				\item \textbf{Maturity:} continuous improvement and metrics
				\item \textbf{Security is continuous risk management}
			\end{itemize}
		\end{block}
	\end{frame}
	
	\begin{frame}{Thank You}
		\begin{center}
			{\Huge \textbf{Good Luck on Your CRISC Exam!}}
		\end{center}
	\end{frame}
	
\end{document}